\documentclass[12pt,a4paper]{article}

% Pacotes de codificação e linguagem
\usepackage[utf8]{inputenc}
\usepackage[T1]{fontenc}
\usepackage[brazil]{babel}
\usepackage{csquotes}

% Pacotes de formatação
\usepackage[margin=2.5cm,top=3cm,bottom=2cm]{geometry}
\usepackage{setspace}
\usepackage{parskip}
\usepackage{titlesec}

% Pacotes gráficos e cores
\usepackage{graphicx}
\usepackage{xcolor}
\usepackage{booktabs}
\usepackage{longtable}
\usepackage{colortbl}
\usepackage{array}

% Pacotes para links e referências
\usepackage[pdftex]{hyperref}
\usepackage{abntex2cite}

% Configuração de cores para links
\hypersetup{
    colorlinks=true,
    linkcolor=blue,
    filecolor=magenta,
    urlcolor=blue,
    citecolor=blue,
    pdftitle={Sistema de Automação para Fábrica de Tijolos Sustentáveis},
    pdfauthor={Felipe Marques},
    pdfsubject={Redes Industriais - Projeto Prático},
}

% Configuração dos títulos das seções
\titleformat{\section}
  {\normalfont\Large\bfseries}{\thesection}{1em}{}
\titleformat{\subsection}
  {\normalfont\large\bfseries}{\thesubsection}{1em}{}
\titleformat{\subsubsection}
  {\normalfont\normalsize\bfseries}{\thesubsubsection}{1em}{}

% Espaçamento entre linhas
\onehalfspacing

% Início do documento
\begin{document}

% Capa
\begin{titlepage}
    \centering
    \vspace*{2cm}

    {\Huge\bfseries Projeto de Rede Industrial\par}
    \vspace{0.5cm}
    {\LARGE\bfseries Automação para Produção Sustentável\par}
    \vspace{1cm}
    {\Large Fábrica de Tijolos Ecológicos EcoBrick\par}

    \vspace{2cm}

    \begin{flushleft}
        \textbf{Disciplina:} Fundamentos de Redes Industriais\\
        \textbf{Professor:} Me. Deivison S. Takatu\\
        \textbf{Data:} Março de 2026
    \end{flushleft}

    \vspace{2cm}

    \begin{flushright}
        \textbf{Aluno:} Felipe Marques\\
        \textbf{Curso:} Análise e Desenvolvimento de Sistemas\\
        \textbf{Instituição:} SENAI
    \end{flushright}

    \vfill

\end{titlepage}

% Sumário
\tableofcontents
\newpage

\section{Introdução e Objetivos}

Este documento apresenta um anteprojeto técnico e financeiro para a modernização da infraestrutura de uma indústria de tijolos ecológicos, focada na implementação de uma rede industrial confiável e escalável. O principal objetivo é substituir a operação manual e isolada por um sistema automatizado integrado, capaz de controlar as variáveis do processo em tempo real e fornecer dados estratégicos para a gestão.

A iniciativa visa não apenas o aumento da produtividade, mas também a garantia da qualidade do produto e a sustentabilidade do negócio, através da redução do consumo de recursos naturais e energia, alinhando-se perfeitamente aos conceitos modernos de automação industrial.

\section{Caracterização do Ambiente Industrial}

\subsection{Descrição do Processo Produtivo}
A unidade fabril em questão é especializada na produção de tijolos modulares de solo-cimento, um material ecológico que dispensa a queima em fornos. O processo é dividido em quatro etapas principais:
\begin{enumerate}
    \item \textbf{Preparação da Matéria-Prima:} Dosagem e mistura homogênea de solo, cimento e água.
    \item \textbf{Prensagem Hidráulica:} A mistura é compactada em prensas hidráulicas para formar os tijolos.
    \item \textbf{Cura Úmida:} Os tijolos são submetidos a um processo de irrigação controlada por aspersores para ganho de resistência.
    \item \textbf{Secagem e Expedição:} Período de repouso em pátio coberto até o transporte.
\end{enumerate}

\subsection{Condições e Desafios Operacionais}
A implementação de uma rede industrial neste ambiente exige a consideração de fatores críticos:
\begin{itemize}
    \item \textbf{Ambiente Agressivo:} Presença constante de poeira fina (solo seco) e alta umidade (setor de cura).
    \item \textbf{Interferência Eletromagnética:} Operação de prensas com motores de alta potência, que geram ruído elétrico significativo.
    \item \textbf{Necessidade de Sincronismo:} A linha de produção requer que as etapas sejam executadas de forma coordenada para evitar gargalos.
    \item \textbf{Monitoramento Remoto:} A área de secagem é extensa e necessita de supervisão à distância.
\end{itemize}

\section{Especificação Técnica e Orçamentária}

Para atender aos requisitos de robustez e confiabilidade, foi realizado um levantamento de mercado com foco em equipamentos com classificação industrial.

\subsection{Unidades de Controle (CLPs)}
Adotou-se uma arquitetura de controle distribuído, com CLPs dedicados para cada etapa do processo, comunicando-se via protocolo industrial.

\begin{longtable}{@{}p{3.5cm}|p{5cm}|c|c@{}}
\caption{Relação de Controladores Lógicos Programáveis} \\
\toprule
\textbf{Modelo / Fabricante} & \textbf{Características Técnicas} & \textbf{Qtd.} & \textbf{Valor Unit. (R\$)} \\
\midrule
\endfirsthead
\toprule
\textbf{Modelo / Fabricante} & \textbf{Características Técnicas} & \textbf{Qtd.} & \textbf{Valor Unit. (R\$)} \\
\midrule
\endhead
\bottomrule
\endfoot
WEG TPW04 & CLP compacto, 24 entradas digitais, 16 saídas a transistor, 2 entradas analógicas, comunicação CANopen e Ethernet. & 4 & 3.200,00 \\
\midrule
WEG TPW04 & Módulo de Expansão de Entradas Analógicas (4-20mA), 8 canais. & 3 & 950,00 \\
\midrule
WEG TPW04 & Módulo de Expansão de Saídas Digitais, 16 canais. & 2 & 720,00 \\
\bottomrule
\end{longtable}
\textbf{Justificativa:} A linha TPW04 da WEG foi escolhida por sua ampla disponibilidade no mercado nacional, suporte técnico local e robustez, sendo ideal para o ambiente fabril brasileiro.

\subsection{Infraestrutura de Comutação (Switches)}
A topologia de rede em anel foi escolhida para garantir alta disponibilidade, utilizando switches gerenciáveis.

\begin{longtable}{@{}p{3.5cm}|p{5cm}|c|c@{}}
\caption{Switches Industriais para Ambiente de Chão de Fábrica} \\
\toprule
\textbf{Modelo / Fabricante} & \textbf{Características Técnicas} & \textbf{Qtd.} & \textbf{Valor Unit. (R\$)} \\
\midrule
\endfirsthead
\toprule
\textbf{Modelo / Fabricante} & \textbf{Características Técnicas} & \textbf{Qtd.} & \textbf{Valor Unit. (R\$)} \\
\midrule
\endhead
\bottomrule
\endfoot
Hirschmann (Belden) & Switch gerenciável, 8 portas RJ45 10/100, 2 portas combo SFP, redundância de anel (MRP), alimentação 24Vcc, grau de proteção IP30. & 5 & 3.100,00 \\
\midrule
Hirschmann (Belden) & Módulo SFP para fibra multimodo, distância até 2km. & 10 & 380,00 \\
\bottomrule
\end{longtable}
\textbf{Justificativa:} Switches gerenciáveis Hirschmann são referência em redes industriais críticas, oferecendo recursos de redundância (MRP) que garantem o funcionamento da rede mesmo em caso de falha de um cabo ou equipamento.

\subsection{Instrumentação (Sensores e Atuadores)}
Os dispositivos de campo são responsáveis pela coleta de dados e execução das ações mecânicas.

\begin{longtable}{@{}p{3.5cm}|p{5cm}|c|c@{}}
\caption{Dispositivos de Campo} \\
\toprule
\textbf{Tipo / Fabricante} & \textbf{Características Técnicas} & \textbf{Qtd.} & \textbf{Valor Unit. (R\$)} \\
\midrule
\endfirsthead
\toprule
\textbf{Tipo / Fabricante} & \textbf{Características Técnicas} & \textbf{Qtd.} & \textbf{Valor Unit. (R\$)} \\
\midrule
\endhead
\bottomrule
\endfoot
Sensor Umidade/Solo & Sonda de umidade para materiais granulares, saída 4-20mA, corpo em aço inox, IP68. & 8 & 1.150,00 \\
\midrule
Transmissor Pressão (Novus) & Faixa 0-200 bar, saída 4-20mA, para sistema hidráulico da prensa, IP66. & 5 & 620,00 \\
\midrule
Sensor Indutivo (IFM) & Chave indutiva M18, alcance 8mm, PNP, para fim de curso de cilindros, IP67. & 30 & 180,00 \\
\midrule
Válvula Solenoide (ASCO) & Válvula direcional 5/2 vias, 24Vcc, para comando de cilindros pneumáticos. & 15 & 290,00 \\
\bottomrule
\end{longtable}

\subsection{Supervisão e Controle (IHM/SCADA)}
Para a interface com o operador e supervisão geral da planta.

\begin{longtable}{@{}p{3.5cm}|p{5cm}|c|c@{}}
\caption{Equipamentos de Supervisão} \\
\toprule
\textbf{Modelo / Fabricante} & \textbf{Características Técnicas} & \textbf{Qtd.} & \textbf{Valor Unit. (R\$)} \\
\midrule
\endfirsthead
\toprule
\textbf{Modelo / Fabricante} & \textbf{Características Técnicas} & \textbf{Qtd.} & \textbf{Valor Unit. (R\$)} \\
\midrule
\endhead
\bottomrule
\endfoot
IHM Delta & Tela touch de 7", resolução 800x480, 2 portas Ethernet, IP65. & 4 & 2.800,00 \\
\midrule
Elipse Software & Elipse SCADA Professional, 2500 tags, driver para CLPs WEG, cliente web incluso. & 1 & 15.900,00 \\
\bottomrule
\end{longtable}

\subsection{Materiais de Rede e Infraestrutura}

\begin{longtable}{@{}p{3.5cm}|p{5cm}|c|c@{}}
\caption{Materiais para Infraestrutura de Rede} \\
\toprule
\textbf{Item} & \textbf{Especificações Técnicas} & \textbf{Qtd.} & \textbf{Valor Total (R\$)} \\
\midrule
\endfirsthead
\toprule
\textbf{Item} & \textbf{Especificações Técnicas} & \textbf{Qtd.} & \textbf{Valor Total (R\$)} \\
\midrule
\endhead
\bottomrule
\endfoot
Cabo de Rede & Cabo Flexível Industrial, SF/UTP, Cat6, antichama, rolo com 305m. & 3 & 1.890,00 (por rolo) \\
\midrule
Cabo de Fibra Óptica & Cabo para externo, armado, 4 fibras multimodo OM3, metro. & 500 & 7,50 (por metro) \\
\midrule
Conector RJ45 & Conector RJ45 metálico, próprio para cabo flexível, blindado. & 150 & 18,00 (unidade) \\
\midrule
Painel Metálico & Painel de aço com visor, 19", 12U, com ventilação e filtro, IP55. & 2 & 1.400,00 \\
\bottomrule
\end{longtable}

\subsection{Resumo do Investimento}

\begin{table}[h]
\centering
\begin{tabular}{@{}p{6cm}>{\raggedleft\arraybackslash}p{4cm}>{\raggedleft\arraybackslash}p{3cm}@{}}
\toprule
\textbf{Categoria} & \textbf{Valor Total (R\$)} & \textbf{(\%)} \\
\midrule
Controladores (CLPs e Módulos) & 16.430,00 & 15,3\% \\
\midrule
Switches e Módulos & 19.300,00 & 17,9\% \\
\midrule
Instrumentação (Sensores/Atuadores) & 22.390,00 & 20,8\% \\
\midrule
Supervisão (IHM/SCADA) & 27.100,00 & 25,2\% \\
\midrule
Infraestrutura e Cabeados & 22.560,00 & 20,8\% \\
\midrule
\textbf{TOTAL GERAL} & \textbf{R\$ 107.780,00} & \textbf{100\%} \\
\bottomrule
\end{tabular}
\caption{Resumo do orçamento consolidado do projeto.}
\end{table}

\section{Avaliação de Adequação e Expansibilidade}

\subsection{Conformidade Técnica com o Ambiente}
A seleção dos equipamentos considerou a severidade do ambiente fabril:
\begin{itemize}
    \item \textbf{Proteção Mecânica:} Todos os componentes de campo possuem grau de proteção mínimo IP65/67, garantindo estanqueidade contra poeira e jatos d'água.
    \item \textbf{Imunidade a Ruídos:} Utilização de cabos blindados (SF/UTP) e conectores metálicos para garantir a integridade dos dados próximo a motores e prensas.
    \item \textbf{Redundância:} A topologia em anel com protocolo MRP assegura que a comunicação seja restabelecida em milissegundos em caso de rompimento de um cabo.
\end{itemize}

\subsection{Perspectivas de Crescimento}
A solução foi planejada para absorver futuras demandas sem a necessidade de substituição de equipamentos principais:
\begin{itemize}
    \item \textbf{Escalabilidade da Rede:} Os switches possuem portas SFP livres, permitindo a expansão da rede para novas áreas ou edifícios utilizando fibra óptica.
    \item \textbf{Capacidade dos CLPs:} Os controladores WEG TPW04 possuem slots para expansão, podendo receber novos módulos de I/O para a adição de sensores.
    \item \textbf{Integração com o Nível Corporativo:} O software Elipse SCADA possui drivers OPC UA e conectores SQL nativos, facilitando a futura exportação de dados de produção (OEE, refugo, consumo) para o sistema ERP da empresa, viabilizando a Indústria 4.0.
\end{itemize}

\section{Estudo de Viabilidade e Análise de Riscos}

\subsection{Análise de Custo-Benefício}
A automação trará ganhos quantificáveis em diversas frentes, conforme análise preliminar:

\begin{table}[h]
\centering
\begin{tabular}{@{}p{5cm}p{5cm}>{\raggedleft\arraybackslash}p{3cm}@{}}
\toprule
\textbf{Indicador de Melhoria} & \textbf{Descrição} & \textbf{Redução/Ganho Anual (R\$)} \\
\midrule
Redução de Perdas & Controle preciso da umidade na cura, diminuindo trincas e quebras. & 35.000,00 \\
\midrule
Eficiência Energética & Desligamento automático de prensas e bombas em horários de inatividade. & 12.000,00 \\
\midrule
Aumento de Produtividade & Redução do tempo de setup entre lotes e operação contínua otimizada. & 60.000,00 \\
\midrule
\textbf{Benefício Total Estimado} & \multicolumn{2}{r}{\textbf{R\$ 107.000,00}} \\
\bottomrule
\end{tabular}
\caption{Projeção de benefícios financeiros anuais.}
\end{table}

\subsection{Retorno sobre o Investimento (ROI)}
Com base nos benefícios projetados e no investimento total:
\begin{itemize}
    \item \textbf{Investimento Total:} R\$ 107.780,00
    \item \textbf{Benefício Líquido Anual Projetado:} R\$ 107.000,00
    \item \textbf{Payback Simples:} Aproximadamente 12 meses.
    \item \textbf{ROI (Ano 1):} 99,3\%
\end{itemize}

\subsection{Análise de Riscos}
A implementação de um projeto desta natureza envolve riscos que precisam ser gerenciados.

\begin{table}[h]
\centering
\begin{tabular}{@{}p{3.5cm}p{2.5cm}p{2.5cm}p{4.5cm}@{}}
\toprule
\textbf{Risco Identificado} & \textbf{Probabilidade} & \textbf{Impacto} & \textbf{Plano de Ação/Mitigação} \\
\midrule
Atraso na importação de componentes & Média & Alto & Priorizar fornecedores nacionais (WEG, Delta) e manter estoque mínimo de peças críticas. \\
\midrule
Interferência na rede PROFINET & Média & Médio & Seguir rigorosamente as normas de instalação (aterramento, separação de cabos) e utilizar equipamentos certificados. \\
\midrule
Resistência da equipe operacional & Alta & Baixo & Realizar treinamentos hands-on com a equipe de manutenção e operadores antes do go-live. \\
\bottomrule
\end{tabular}
\caption{Matriz de riscos e estratégias de mitigação.}
\end{table}

\section{Arquitetura da Solução Proposta}

\subsection{Topologia Lógica e Física}
A arquitetura proposta segue o modelo de pirâmide de automação, segregando os níveis de campo, controle e supervisão.

\begin{enumerate}
    \item \textbf{Nível de Campo (Dispositivos):} Sensores de umidade/pressão e válvulas solenoides conectados diretamente aos CLPs via cabos elétricos (4-20mA e 24Vcc).
    \item \textbf{Nível de Controle (PLC):} Quatro CLPs WEG TPW04 distribuídos estrategicamente (um para dosagem, dois para prensagem, um para cura). Eles se comunicam entre si via rede PROFINET em anel.
    \item \textbf{Nível de Supervisão (SCADA/IHM):}
        \begin{itemize}
            \item Um servidor com Elipse SCADA coleta dados de todos os CLPs para monitoramento central, geração de relatórios históricos e alarmes.
            \item Quatro IHMs Delta instaladas em pontos-chave (prensa, expedição, etc.) permitem a operação local da máquina.
        \end{itemize}
    \item \textbf{Conexão Corporativa:} Um switch de borda (firewall) separa a rede industrial (192.168.1.x) da rede corporativa (10.0.0.x), garantindo segurança. Futuramente, um gateway OPC UA fará a ponte de dados entre o SCADA e o ERP.
\end{enumerate}

\subsection{Representação Visual (Descrição do Diagrama)}
O diagrama proposto ilustra uma topologia em estrela estendida. No centro, um switch industrial Hirschmann faz o papel de "core" da rede industrial. Conectados a ele, em anel, estão os switches distribuídos nos três painéis elétricos da fábrica. Cada um desses switches conecta os CLPs e IHMs locais. O servidor SCADA está conectado diretamente ao switch core, assim como o firewall que faz a ponte com o restante da empresa. As linhas pontilhadas representam o backbone de fibra óptica entre os painéis, garantindo alta velocidade e imunidade a ruídos.

\section{Conclusão e Recomendações}
A análise realizada confirma a viabilidade técnica e econômica da implantação de uma rede industrial na fábrica de tijolos ecológicos. A proposta não apenas resolve os desafios atuais de controle e monitoramento, mas também estabelece uma base sólida para a transformação digital da empresa.

O investimento de R\$ 107.780,00 apresenta um retorno atraente, com payback em torno de 12 meses, impulsionado pela redução de perdas e ganhos de eficiência. Recomenda-se a adoção da arquitetura proposta, priorizando a contratação de mão de obra especializada para a instalação e configuração, garantindo que os rigorosos padrões de redes industriais sejam atendidos e o sistema opere com o máximo de determinismo e confiabilidade.

\begin{thebibliography}{9}

\bibitem{weg} WEG Equipamentos Elétricos S.A. \textit{Manual do Programador CLP TPW04}. Disponível em: \url{https://www.weg.net}. Acesso em: 10 mar. 2026.

\bibitem{belden} Belden Inc. \textit{Hirschmann Industrial Ethernet Switches Datasheet}. Disponível em: \url{https://www.belden.com}. Acesso em: 10 mar. 2026.

\bibitem{elipse} Elipse Software. \textit{Manual do Usuário Elipse SCADA}. Disponível em: \url{https://www.elipse.com.br}. Acesso em: 11 mar. 2026.

\bibitem{lara} LARA, Carla Eduarda Orlando de Moraes de. \textit{Automação e controle industrial}. Curitiba: Contentus, 2021.

\bibitem{lugli} LUGLI, Alexandre Baratella; SANTOS, Max Mauro Dias. \textit{Sistemas Fieldbus para automação industrial}. São Paulo: Érica, 2009.

\bibitem{moraes} MORAES, Cícero Couto de; CASTRUCCI, Plinio de Lauro. \textit{Engenharia de automação industrial}. São Paulo: LTC, 2007.

\bibitem{takatu} TAKATU, Deivison S. \textit{Fundamentos de Redes Industriais}. Material didático, SENAI, 2026.

\end{thebibliography}

\end{document}